% Pursuing Stacks (cotes 134-2 to 134-8): archive cote 134-4, site page 12 (PDF page 13), typescript page 100.
% typescript. Draft compared: ps3.tex:592-650.
% AI-transcribed by Claude Opus 5.5 (claude-opus-5-5, adaptive thinking, effort high), Message Batches API, batch ps1 (msgbatch_01Rv5hCaNTtUt4JHx1yqoHPL), 2026-11-12; not yet checked by eye.
\scanpage{134-4/12}
\mspage{(100)}
vision of \hand{$\Delta$}$_1$,\xout{withxxixx} which can also be interpreted as an\xout{x} amalgamated
sum of \hand{$\Delta$}$_1$ with itself, namely
\[ I \ = \ \diagram{\begin{tikzcd}[row sep=small, column sep=small] e_0 \arrow[dr] & \\ & e_2 \\ e_1 \arrow[ur] & \end{tikzcd}} \qquad , \]
an $\underline{I}$-homotopy from f to g is essentially a triple (k,u,v), with
k: $X \hand{$\to$} Y$ and u: $f \hand{$\to$} k$ and v: $g \hand{$\to$} k$ maps in \uline{Hom}(X,Y); this time, the
relation of I-homotopy is symmetric, but by no means transitive.

Returning to general M, it is customary to introduce the equivalen-%
ce relation in Hom(X,Y) generated by the relation of $\underline{I}$-homotopy - we'll
say that f and g are \uline{$\underline{I}$-homotopic in the wide sense}, and we'll write
\[ f \underset{\underline{I}}{\hand{$\sim$}} g \qquad , \]
if they are equivalent with respect to this relation. As seen from the \xout{fx}
\deleted{first two} examples above (where M = (Cat))\hand{,} this relation still depends on
the choice of the homotopy interval $\underline{I}$ . Let's first look at what we can
do for fixed I, and then how what we do depends on I.

If f and g are I-homotopic, then so are their compositions with
any $Y \hand{$\to$} Z$ or $T \hand{$\to$} X$. This implies that the relation $\underset{\underline{I}}{\hand{$\sim$}}$ of $\underline{I}$-homotopy
in the wide sense is compatible with compositions. If we denote by
\[ \mathrm{Hom}_{\deleted{\illegible{1}}}(X,Y)_{\underline{I}} \]
the quotient set of Hom(X,Y) by the equivalence relation of $\underline{I}$-homotopy in
the wide sense, we get composition between the sets Hom(\deleted{\illegible{1}}X,Y)$_{\underline{I}}$ , and
hence a structure of a category $M_I$ having same objects as M, and where
maps from X to Y are $\underline{I}$-homotopy classes (in the wide sense - this will
be understood henceforth when speaking of homotopy classes) of maps from
X to Y in M. Two objects X,Y of M i.e. of $M_I$ which are isomorphic as
objects of $M_I$ will be called \uline{$\underline{I}$-homotopic}. This means also that we can
find maps in M \inserted{(so-called \uline{$\underline{I}$-homotopisms} - namely $M_{\underline{I}}$-\unsure{isomorphims|isomorphisms})}
\[ f\colon X \hand{$\to$} Y \quad , \quad g\colon Y \hand{$\to$} X \]
such that \xout{thexcompx} we get $\underline{I}$-homotopy relations in the wide sense
\[ gf \underset{\underline{I}}{\hand{$\sim$}} \mathrm{id}_X \quad , \quad fg \underset{\underline{I}}{\hand{$\sim$}} \mathrm{id}_Y \ . \]
We'll say X is \uline{$\underline{I}$-contractible} if X is $\underline{I}$-homotopic to the final object $e_M$=
e of M (which visibly is also a final object of $M_I$), i.e.\xout{thisx} if X is a
final object of $M_I$ . In terms of M, this means that there exists a
section f of X over e, such that $fp_X$ is $\underline{I}$-homotopic in the wide sense
to $\mathrm{id}_X$ (where $p_X$ is the unique map $X \hand{$\to$} e$). In fact, if there is such a
section f, any other section will do too.

From the homotopy lemma a) it follows that the canonical functor
\xout{MxxxHx}
\margin{12}
